\begin{phasebox}
\textbf{هدف:} مشخص‌کردن شکل پیام‌ها قبل از نوشتن کد دو طرف ارتباط.\\
\textbf{پیش‌نیاز:} یک شناسه آزمایشی برای دستگاه.\\
\textbf{معیار پایان:} نام موضوع و نمونه پیام داده، وضعیت، فرمان و پاسخ فرمان تعریف شوند.
\end{phasebox}

\section{قرارداد چیست؟}
قرارداد در این پروژه یک فایل ساده است که مشخص می‌کند برد و برنامه پایتون چه پیام‌هایی را با چه نام و چه فیلدهایی جابه‌جا می‌کنند.

\section{موارد لازم}
برای ارتباط
\lr{MQTT}
نام
\lr{Topic}
فرستنده، گیرنده، سطح
\lr{QoS}
و وضعیت نگهداری پیام مشخص می‌شود. برای پیام
\lr{JSON}
نیز نام فیلدها، نوع مقدار، واحد، زمان و اجباری‌بودن هر فیلد ثبت می‌شود.

\section{محل نگهداری}
همه تعریف‌ها و نمونه پیام‌ها مستقیماً در پوشه
\lr{\texttt{code/contracts}}
قرار می‌گیرند. در نسخه اول، فایل متنی و نمونه
\lr{JSON}
کافی است.

\section{آزمایش مستقل}
پیام نمونه با ابزار خط فرمان منتشر و دریافت می‌شود. سپس همان پیام در تست برنامه پایتون و برنامه برد استفاده خواهد شد.

\section{قاعده تغییر}
هر تغییری که کد قبلی را خراب می‌کند باید با شماره نسخه جدید قرارداد انجام شود.
